\scanpage{541}
\mspage{276}

\unsure{(12)} où $\Sigma$ \inserted{$\subset\mathbb C$} est \unsure{la sous-extension} cyclotomique maximale de $\overline{\mathbb Q}_0$.

\deleted{\unsure{Le}} On récupère les $M_{1,1}[n']$, pour $n'\mid n$, \inserted{$n'\ne 1,2$,} à partir de $M_{1,1}[n]$ avec l'action de $\Gamma_n$ dessus, par
\[
(13)\qquad M_{1,1}[n']\simeq M_{1,1}[n]\wedge_{\Gamma_n}\Gamma_{n'}\simeq M_{1,1}[n]/\unsure{K_{n,n'}}
\]
\[
\text{sur }\operatorname{Spec}\mathbb Z[1/n]
\]
où
\[
(14)\qquad K_{nn'}=\operatorname{Ker}(\Gamma_n\longrightarrow\Gamma_{n'}).
\]

\margin{\uline{NB} $K_{n,n'}$ opère \uline{librement} sur le schéma $M_{1,1}[n]$ (si $n'\mid n$, $n'\ne 1,2$) --- mais pas bien sûr si $n'=1$ ou $2$.}

Si on applique cependant cette formule dans le cas non \unsure{licite} $n'=1$ ou $2$, on trouve pour le cas $n'=1$, non $M_{1,1}$ mais $\widetilde M_{1,1}$ (schéma modulaire grossier), et pour $n'=2$, non $M_{1,1}[2]$ (qui n'est pas non plus un schéma) mais ce qu'on avait appelé au \S\ précédent $M_{1,1}[2]'$.

On peut \deleted{\unsure{déterminer}} \inserted{expliciter} \inserted{$M_{1,1}$ et} $\widetilde M_{1,1}$ sur $\operatorname{Spec}\mathbb Z$ tout entier, \inserted{en \unsure{termes} des $M_{1,1}[n]$,} en les décrivant \unsure{au} \inserted{dessus de} deux ouverts $\operatorname{Spec}(\mathbb Z[1/n])$, $\operatorname{Spec}(\mathbb Z[1/n'])$ qui recouvrent $\operatorname{Spec}(\mathbb Z)$ (i.e. pour $(n,n')=1$) par (9) et (10), \inserted{(en y prenant d'abord $n$, puis $n'$)} et en ``recollant'' \deleted{\unsure{des|de}} au dessus de $\operatorname{Spec}(\mathbb Z[\frac{1}{nn'}])$ par ces $\overline{\mathrm m}$s formules, appliquées à $nn'$\ldots

\scanpage{542}

On a d'ailleurs (pour mémoire)
\[
\text{\uline{Th}}\qquad \widetilde M_{1,1}\simeq\mathbb E^1_{\mathbb Z}\qquad(\text{sur }\operatorname{Spec}\mathbb Z)
\]

On peut épingler un tel isomorphisme (\unsure{défini|définit} à priori modulo le groupe affine entier $X\longmapsto\pm X+n$, $n\in\mathbb Z$) en \unsure{notant} qu'il y a dans $\widetilde M_{1,1}$ deux sections privilégiées sur $\mathbb Z$, dont les valeurs au pt générique correspondent aux deux classes d'isom des courbes elliptiques en car. $0$ (sur $\mathbb C$, p.ex) qui ont des autom. $\ne\mathrm{id},-\mathrm{id}$ --- les deux groupes \unsure{d'automs} \unsure{qu'on obtient} sont d'ailleurs $\mathbb Z/4\mathbb Z$ et $\mathbb Z/6\mathbb Z$ (très facile p.ex. par voie \unsure{transcend.te|transcendante}). Ce sont ce qu'on appelle sauf erreur les courbes elliptiques ``anharmoniques''. Sur un $M_{1,1}[n]$, le groupe des autom. rat.\,$/k$ d'une courbe elliptique décrite par un point $x$ de $M_{1,1}[n](k)$ s'identifie canoniquement au stabilisateur de $x$ dans $\Gamma_n$ :
\[
(15)\qquad\operatorname{Aut}(E_x)\simeq(\Gamma_n)_x
\]

\scanpage{543}
\mspage{277}

--- bien sûr, on a des énoncés \unsure{idoines|identiques} sur un schéma de base quelconque, pas néc. un corps. Cela montre déjà que les \unsure{automs} des courbes elliptiques sont liés de façon essentielle à la \uline{ramification} de $\Gamma_n$ opérant sur $\unsure{M_{1,1}[n]|M_{1,n}[n]}$. On trouve ainsi :

\uline{Proposition} (Sur $\operatorname{Spec}\mathbb Q$) Il y a exactement deux orbites de $\Gamma_n$ opérant sur $M_{1,1}[n]$ qui sont ``\deleted{singulières}\inserted{critiques}'', et qui correspondent à des stabilisateurs isomorphes resp. à $\mathbb Z/4\mathbb Z$ et $\mathbb Z/6\mathbb Z$, qui sont les groupes des \unsure{automs} des courbes elliptiques (``anharmoniques'') correspondantes.

\margin{$n\ne 1,2$}

On \unsure{notera|note} que $\Gamma_n$ \uline{n'opère pas} \unsure{fidèlt|fidèlement} sur $M_{1,1}[n]$ --- le noyau de cette opération est le s-groupe central $\{\pm1\}$, image du groupe correspondant de $\operatorname{Gl}(2,\mathbb Z)$ --- il correspond au groupe des automorphismes ``universels'' $\pm\mathrm{id}_E$ des courbes elliptiques $E$. Le groupe quotient
\[
(16)\qquad \Gamma'_n=\Gamma_n/\{\pm1\}
\]
\deleted{\unsure{On|sur}} opère \unsure{fidèlt|fidèlement}, et on peut regarder le ``topos mixte'' $(M_{1,1}[n],\Gamma'_n)$ --- on trouve

\scanpage{544}
qu'il ne \uline{dépend pas} (à \unsure{équi.|équiv.} can. près) du \uline{choix de $n$} [à cela près que l'ouvert de $\operatorname{Spec}\mathbb Z$ sur lequel ``il a un sens'' dépend de $n$\dots] -- et. en fait, on a :

\margin{donc à \unsure{condition} de se placer au-dessus d'un schéma de base (tel $\operatorname{Spec}\mathbb Z[\frac{1}{nn'}]$) où $n,n'$ sont devenus inversibles\dots}

\[
(17)\qquad (M_{1,1}[n],\Gamma'_n)\simeq M_{1,1}'
\]
\[
(\text{au dessus de }\operatorname{Spec}\mathbb Z[1/n])
\]
où le \deleted{topos} \inserted{multiplicité schématique} modulaire $M_{1,1}'$ est décrit, en termes des courbes elliptiques ``\unsure{définies} mod symétrie'', comme au \S\ précédent.

\uline{Remarque.} La proposition précédente semble dépendre de $n$ -- mais on voit a priori \deleted{que} grâce au fait que les $\Gamma_{n,n'}$ \inserted{($n'\mid n$, $n'\ne1,2$)} (\unsure{opérant|opèrent} librement) que si elle est valable pour \uline{un} $n$, elle est valable pour tous.

\deleted{\unsure{Cor.}} Notons \deleted{que} aussi

\uline{Corollaire de la proposition.} \inserted{(car. 0)} Les stabilisateurs \uline{dans $\Gamma'_n$} (le groupe qui opère \uline{\unsure{fidèlement}} sur $M_{1,1}[n]$) des points des \unsure{deux} orbites critiques sont resp. $\mathbb Z/2\mathbb Z$ et $\mathbb Z/3\mathbb Z$. \deleted{\illegible{1}}

\margin{valable pour $n\unsure{\ne|=}2$, \unsure{cf|cf.} plus bas}

\scanpage{545}
\mspage{278}
On a supposé ici que $n\ne1,2$ -- mais rien ne nous empêche de regarder aussi \inserted{sur}
\[
(18)\qquad M_{1,1}[2]'\overset{\mathrm{dfn}}{=}M_{1,1}[2m]/\Gamma_{2m,2}
=M_{1,1}[2m]/\unsure{\Gamma_{2m,2}|\Gamma'_{2m,2}}
\]
(où $m$ est un entier quelconque, $\ge2$)
l'action de $\Gamma_2\simeq\operatorname{Gl}(2,\mathbb Z/2\mathbb Z)\simeq\operatorname{Sl}(2,\mathbb Z/2\mathbb Z)\simeq\unsure{\mathfrak S_3|\mathcal S_3}\simeq\Gamma_{2m}/\Gamma_{2m,2}$. \deleted{\unsure{Alors que}} Il n'est plus vrai, \inserted{certes,} que $\Gamma_{2m,2}$ opère \unsure{librement} sur $M_{1,1}[2m]$ -- car $\Gamma_{2m,2}$ contient l'élément $-1\in\Gamma_{2m}$ correspondant à la symétrie universelle $-\operatorname{id}_E$ -- i.e. \inserted{\unsure{à}} l'élément \deleted{\unsure{de}} $-1$ de $\operatorname{Gl}(2,\mathbb Z)$. Mais le résultat général de rigidité rappelé au début implique que \deleted{$\Gamma'_{2m,2}=\Gamma_{2m,2}/\pm1$}
\[
(19)\qquad\Gamma'_{2m,2}\overset{\mathrm{dfn}}{=}\Gamma_{2m,2}'/\pm1
\simeq\operatorname{Ker}(\Gamma'_{2m}\longrightarrow\Gamma'_2\simeq\unsure{\mathfrak S_3|\mathcal S_3})
\]
opère \deleted{\illegible{1}} \inserted{\unsure{encore}} \uline{librement} sur $M_{1,1}[2m]$ -- et cela implique que le cor. de la prop. est valable encore pour l'action de $\Gamma'_2\,(\unsure{\xleftarrow{\sim}|\simeq}\Gamma_2\simeq\unsure{\mathfrak S_3|\mathcal S_3})$ sur $M_{1,1}[2]'$.

\uline{Digression terminologique.} $M_{1,1}$ \inserted{${}=M_{1,1}[1]$} et les

\scanpage{546}
\mspage{278}
\figure{[Le numéro 278 est entouré d'un ovale; une flèche horizontale pointe vers le numéro 278' à droite.]}
\mspage{278'}
$M_{1,1}[n]$ ($n\in\mathbb N^*$) sont \unsure{définis} \unsure{a|à} priori comme des \uline{multiplicités} schématiques modulaires, i.e. a) comme des topos loc. annelés loc. isom. au topos étale d'un schéma (en l'occurrence, de type fini sur $\operatorname{Spec}\mathbb Z$, ou sur $\mathbb Q$ si on travaille en car. 0) et b) \unsure{définis}, à équivalence près définie elle-$\overline{\mathrm m}$ à \unsure{isom|isom.} unique près, comme ``représentant'' les \deleted{\illegible{1}} ``foncteurs'' correspondants
\[
S\longmapsto
\begin{gathered}
\text{catégorie des courbes elliptiques}\\
\text{relatives sur }S\text{, avec rigidification}\\
\text{de Jacobi d'échelon }n\text{ i.e.}\\
\text{iso.}\quad(\mathbb Z/n\mathbb Z)_S\simeq{}_nE\\
={}\unsure{\mathcal{E}ll_n}(S)
\end{gathered}
\]
donc par construction on a \inserted{pour \unsure{tt|$\forall$} schéma $S$} une équivalence
\[
\unsure{\mathcal{E}ll_n}(S)\unsure{\simeq|\cong}
\text{\uline{$\operatorname{Hom}$}}_{\substack{\text{topos loc.}\\\text{annelés}\\\text{i.e. hom de}\\\text{multiplicités}}}
(S,M_{1,1}[n])
\]
Pour que les catégories $\unsure{\mathcal{E}ll_n}(S)$ soient \uline{discrètes}, \deleted{\illegible{1}} ou encore $M_{1,1}[n]$ soit un schéma, ou encore que les autom. d'une

\scanpage{547}
\mspage{279}

courbe ell. \deleted{\unsure{respectant|respectent}} une rigid. de Jacobi d'échelon $n$ \unsure{soient} l'identité, il \unsure{f et s.} que $n\geq 3$ i.e. $n\neq 1,2$. Donc $M_{1,1}[2]$ n'\uline{est pas} un schéma -- mais $M'_{1,1}[2]$ en est un, et $M_{1,1}[2]\to M_{1,1}[2]'$ est un \deleted{\illegible{1}} isom local (décrit entièrement par une $\mathbb Z/2\mathbb Z$-gerbe sur $M_{1,1}[2]'$). \deleted{\unsure{On considère bien sûr}} \inserted{\unsure{Il} faut quand $\overline{\mathrm m}$ prendre la} peine de définir
\diagram{\[
\tag{20}
\begin{tikzcd}[ampersand replacement=\&,column sep=large,row sep=large]
M_{1,1}[2] \arrow[r] \arrow[d,phantom,"{\simeq}"{description,rotate=90},"{\text{``théor.''}}"{description,xshift=3em}] \& M_{1,1}[2]' \arrow[d,phantom,"{\simeq}"{description,rotate=90},"{\text{``\unsure{déf\(^{\mathrm n}\)|déf}''}}"{description,xshift=3em}] \\
(M_{1,1}[2m],\Gamma_{2m,2}) \& M_{1,1}[2m]/\Gamma'_{2m,2} \arrow[d,phantom,"{\Vert}"{description}] \\
\& M_{1,1}\unsure{(2m)|[2m]}/\Gamma_{2m,2}
\end{tikzcd}
\]}
comme l'homomorphisme canonique du ``passage au quotient grossier'' -- qui a un sens chaque fois qu'un groupe discret (disons)\unsure{\S} $G$ opère sur un schéma (disons) $X$, comme un morphisme
\[
\tag{21}
(X,G)\longrightarrow X/G=Y
\]
dont le foncteur \unsure{image} inverse (faisceaux sur $Y=X/G\longrightarrow G$-faisceaux sur $X$) est \unsure{évidente|évident}. Ici on a un sous-groupe inv. \unsure{$\mathfrak z$|$\mathfrak Z$} de $G$ qui opère \unsure{trivialt|trivialement} sur $X$, tel que $G/\unsure{\mathfrak z|\mathfrak Z}$ opère librement, ce qui signifie que le morphisme

\scanpage{548}

(21) se factorise en
\[
\tag{22}
(X,G)\longrightarrow
\left(X,\underbrace{G/\unsure{\mathfrak z|\mathfrak Z}}_{G'}\right)
\overset{\sim}{\longrightarrow}Y=X/G=X/G',
\]
de sorte que le morphisme de multiplicités (21) s'identifie \unsure{ainsi,|ainsi} simplement, à
\diagram{\[
\tag{23}
\begin{tikzcd}[ampersand replacement=\&,column sep=large,row sep=large]
(M_{1,1}[2m],\Gamma_{2m,2}) \arrow[r] \& (M_{1,1}[2m],\Gamma'_{2m,2}) \arrow[d,phantom,"{\simeq}"{description,rotate=90}] \\
\& M_{1,1}[2]'
\end{tikzcd}
\]}

Par contre, si on veut descendre jusqu'à $M_{1,1}[1]'=M_{1,1}'$, en passant au quotient (au sens ``grossier'') de $M_{1,1}[n]$ par $\Gamma'_n$, le morphisme
\diagram{\[
\tag{24}
\begin{tikzcd}[ampersand replacement=\&,column sep=large,row sep=large]
(M_{1,1}[n],\Gamma'_n) \arrow[r] \arrow[d,phantom,"{\Vert}"{description}] \& M_{1,1}[n]/\Gamma'_n \arrow[d,phantom,"{\Vert}"{description}] \\
M_{1,1}' \& \widetilde M_{1,1}
\end{tikzcd}
\]}
n'est plus un ``iso'' i.e. n'est plus une équivalence, car $\Gamma'_n$ n'opère pas librement -- il a \inserted{($\overline{\mathrm m}$ en car. $0$)} de la ramification de degré $2,3$.

[Si cependant on exclut les courbes elliptiques anharmoniques, \unsure{sa|ça} marcherait -- ce n'est ``qu'au voisinage'' de ces courbes elliptiques que (24) n'est pas un iso\ldots

\scanpage{549}
\mspage{280}

En fait, les raisonnements du \S{} précédent établissent un isomorphisme (\unsure{valable} sur \inserted{le} $\mathrm{Spec}\,\unsure{\mathbb Z[1/n]}$ de $\mathbb Z[1/n]$)
\[
(25)\qquad M_{1,1}[2]'\simeq M_{0,4}^{!}\simeq U_{0,3}
\]
(\unsure{valable} sur $\mathrm{Spec}\,\mathbb Z[1/2]$)

\margin{NB \unsure{l'image} de $M_{1,1}[2]'$ dans $\mathrm{Spec}\,\mathbb Z$ est $\mathrm{Spec}\,\mathbb Z[1/\unsure{n|2}]$, tandis que celle de $M_{0,4}^{!}\simeq U_{0,3}$ est $\mathrm{Spec}\,\mathbb Z$ \unsure{tt} entier}

et cet isomorphisme est compatible avec les opérations de \deleted{$\unsure{\Gamma'_2\simeq\Gamma}$}
\[
(26)\qquad \Gamma'_2\simeq\Gamma_2=\operatorname{Gl}(2,\mathbb Z/2\mathbb Z)=\operatorname{Sl}(2,\mathbb Z/2\mathbb Z)
\]
sur $M_{1,1}[2]'$ d'une part, de $\unsure{\mathfrak S_3|\mathcal S_3}$ sur $U_{0,3}$ d'autre \unsure{part|parts}, quand on considère l'isomorphisme
\[
(27)\qquad
\begin{array}{ccc}
\Gamma_2 & \xrightarrow{\sim} & \unsure{\mathfrak S_3|\mathcal S_3}\\
\Vert &&\\[-4pt]
\operatorname{Gl}(2,\mathbb F_2)&&
\end{array}
\]
qui associe, à tout autom de $\mathbb F_2^2$, son action sur les trois éléments non nuls, $\underline{\alpha}=(1,0)$, $\inserted{\underline{\beta}=}(0,1)$, et $\underline{\gamma}=(1,1)=\underline{\alpha}+\underline{\beta}$ (de sorte qu'on a d'ailleurs) \deleted{\illegible{2}}
\[
(28)\qquad \gamma=\alpha\unsure{+|-}\beta,\quad \unsure{\alpha|\beta}=\beta\unsure{+|-}\gamma,\quad \beta=\unsure{\gamma}+\alpha.
\]
Il faudrait expliciter cette compatibilité, en \unsure{considérant} l'action naturelle de $\unsure{\mathfrak S_n|\mathcal S_n}$ sur $M_{0,4}^{!}$, $M_{0,n}^{!}$ étant la multiplicité modulaire définie par

\scanpage{550}

\[
(29)\qquad
\underline{\operatorname{Hom}}_{\text{multiplicités}}(S,M_{0,n}^{!})\simeq
\begin{gathered}
\text{catégorie des}\\
\text{\unsure{courbes|droites} relatives}\\
\text{sur }S,\ \text{loc. isom. à }\mathbb P^1_S\\
\text{(pour top. étale) et munies}\\
\text{de }n\text{ sections mutuellement}\\
\text{\unsure{disjointes} \unsure{numérotées}}\\
\deleted{\illegible{1}}\quad\unsure{s_1|s_0},s_2,\ldots s_n
\end{gathered}
\]

\margin{Si on veut une représentabilité pour $n=0,1,2$, il ne suffit plus de travailler avec \unsure{la topologie étale|les topologies étales} top. fppf p.ex...}

\uline{NB} c'est \inserted{\unsure{co|en}} représentable bel et bien par une multiplicité ssi $n\geq3$, \inserted{i.e. si le champ \unsure{des groupoïdes} $\unsure{S\mapsto|S\to}$ \unsure{est représentable et} ``rigide'' ---} et alors la multiplicité est $=$ un schéma, isomorphe d'ailleurs \unsure{canonique|canoniquement} au \inserted{\unsure{ss-}}schéma
\[
\underline{\operatorname{Mon}}(I_{n-3},U_{0,3})\subset U_{0,3}^{I_{n-3}},
\]
où ici
\[
I_{n-3}=\{i\in\mathbb N\quad3<i\leq n\deleted{\unsure{-3}}\}
\]
--- ainsi \deleted{\illegible{1}}
\[
M_{0,3}^{!}\xrightarrow{\sim}\mathrm{Spec}\,\mathbb Z\quad\text{(schéma final)},\quad\text{et}
\]
\[
M_{0,4}^{!}\xrightarrow{\sim}U_{0,3}=\mathbb P^1_{\mathbb Z}\setminus\{\text{sections }0,1,\infty\}=\mathrm{Spec}\,\mathbb Z[T]\Bigl[\frac{1}{T(T-1)}
\]

Mais on fera attention qu'il y a une opération naturelle de $\unsure{\mathfrak S_n|\mathcal S_n}$ sur $M_{0,n}^{!}$, donc sur $\underline{\operatorname{Mon}}(I_{n-3},U_{0,3})$ --- \deleted{\unsure{dont|donc}} alors que l'on ne voit a priori que l'action de $\unsure{\mathfrak S_3|\mathcal S_3}\times\unsure{\mathfrak S_{n-3}|\mathcal S_{n-3}}$ sur ce dernier: l'isom canonique
\[
(30)\qquad M_{0,n}^{!}\simeq\unsure{\operatorname{Mon}}(I_{n-3},U_{0,3})\subset U_{0,3}^{I_{n-3}}
\]
\[
\text{(où }I_{n-3}=\{i\in\mathbb N\mid\deleted{\illegible{1}}3<i\leq n\}\text{)}
\]
ne tient compte que des opérations naturelles \unsure{des}
